\chapter{Vector and tensor fields}
\label{chap:dev:vector-tensor-fields}

\begin{quote}
Phase metadata. Spec dir:
\passthrough{\lstinline!specs/2026-05-04-phase-13-vector-tensor-fields/!}.
Generalizes \passthrough{\lstinline!Field<T>!} to a rank-2 element
type and adds rank-raising / rank-lowering
\passthrough{\lstinline!gradient!} and
\passthrough{\lstinline!divergence!} overloads, plus the contraction
primitives \passthrough{\lstinline!trace!},
\passthrough{\lstinline!singleContract!}, and
\passthrough{\lstinline!doubleContract!}. Version target:
\passthrough{\lstinline!0.13.0!}.
\end{quote}

\subsection{1. Theory}\label{phase13-theory}

\subsubsection{Rank-raising gradients on continuous
fields}\label{phase13-continuous}

For a smooth scalar field \(\phi: \mathbb{R}^3 \to \mathbb{R}\), the
gradient is a rank-1 covector field
\((\nabla\phi)_j = \partial_j\phi\). For a smooth vector field
\(v: \mathbb{R}^3 \to \mathbb{R}^3\), the natural rank-raising
operator is the Jacobian-style rank-2 tensor field
\((\nabla v)_{ij}\), traditionally defined either as
\((\nabla v)_{ij} = \partial_j v_i\) (the velocity-gradient
convention used in continuum mechanics) or as
\((\nabla v)_{ij} = \partial_i v_j\) (some textbooks). Phase 13
adopts the velocity-gradient form so that the trace recovers the
divergence:
\[
\mathrm{tr}(\nabla v) = (\nabla v)_{ii} = \partial_i v_i = \nabla\cdot v.
\]
Combined with the matched divergence
\((\mathrm{div}\,M)_i = \partial_j M_{ij}\), the composition
\((\mathrm{div}(\nabla v))_i = \partial_j\partial_j v_i\) is the
\(i\)-th Cartesian component of the vector Laplacian — the third
roadmap-acceptance identity.

\subsubsection{Schwarz's theorem and the
curl-of-gradient identity}\label{phase13-schwarz}

For \(\phi \in C^2\), Schwarz's theorem (also: Clairaut's theorem)
gives
\(\partial_i\partial_j\phi = \partial_j\partial_i\phi\). The Hessian
\(H_{ij} := \partial_j\partial_i\phi\) is therefore symmetric, which
is equivalent to the vector identity
\(\nabla\times\nabla\phi \equiv 0\) on simply connected domains:
the antisymmetric part of the Hessian, recoverable as the curl of the
gradient, vanishes identically. Phase 13 verifies this in
rank-raising vocabulary by computing
\passthrough{\lstinline!gradient(gradient(phi))!} directly and
asserting symmetry of the resulting
\passthrough{\lstinline!Field<Mat3d>!} pointwise.

\subsection{2. Math derivation}\label{phase13-math}

\subsubsection{Discrete commutativity of partial-difference
operators}\label{phase13-commute}

Let \(D_a\) denote the discrete first-derivative operator along axis
\(a\) using the
\passthrough{\lstinline!stencil::FirstDerivative<Order>!} weights
\((w[m], \mathit{off}[m])\). On a periodic uniform grid:
\[
(D_a\,\psi)_p \;=\; h_a^{-1} \sum_m w[m]\,\psi_{p + \mathit{off}[m]\,e_a}.
\]
Composing two such operators along (possibly equal) axes \(a\), \(b\):
\[
(D_b\,D_a\,\psi)_p \;=\; (h_a h_b)^{-1} \sum_{m, m'} w[m]\,w[m']\,\psi_{p + \mathit{off}[m]\,e_a + \mathit{off}[m']\,e_b}.
\]
The right-hand side is symmetric in the pair \((a, m) \leftrightarrow
(b, m')\): swapping the two summation pairs leaves the result
unchanged. Therefore \(D_b\,D_a\,\psi = D_a\,D_b\,\psi\) at every
grid point — the discrete partial-difference operators commute on a
periodic uniform grid \emph{exactly} (modulo round-off in the
accumulated sum), not merely at convergence rate. Hessian symmetry is
therefore a round-off-level property of the discrete computation, not
a stencil-order one. The Phase 13 test asserts a max-asymmetry below
\(10^{-12}\).

\subsubsection{Trace = divergence pointwise to
round-off}\label{phase13-trace}

The rank-2 gradient overload computes, at each grid point,
\(M(0, 0) = (D_x v_x)\), \(M(1, 1) = (D_y v_y)\), and
\(M(2, 2) = (D_z v_z)\). The Phase 2 scalar divergence applied to the
same \(v\) computes the same three quantities and sums them in a
single per-point accumulation. The two routes therefore agree
pointwise to round-off — only the order of the floating-point
summation can differ.

\subsubsection{Vector Laplacian via composition: the stride-2
phenomenon, revisited}\label{phase13-stride2}

For the discrete divergence-of-gradient on a vector field:
\[
(\mathrm{div}(\nabla v))_i \;=\; D_j\,M(i, j) \;=\; D_j\,(D_j v_i) \;=\; (D_j)^2 v_i,
\]
i.e., the discrete operator
\((D_j)^2\) acts on \(v_i\) componentwise and the three axis
contributions are summed. Each \((D_j)^2\) is the composition of two
\passthrough{\lstinline!FirstDerivative<2>!} stencils: at offsets
\((-1, 1)\), \((-1, 1)\), the resulting effective stencil reaches
offsets \((-2, 0, +2)\) along axis \(j\) — a \emph{stride-2}
second-derivative stencil. Compare with
\passthrough{\lstinline!diff::laplacian!}, which uses the 3-point
\((-1, 0, +1)\) stencil from
\passthrough{\lstinline!Coefficients<2>!}. Both stencils are
\(\mathcal{O}(h^2)\) accurate but with different leading-error
constants (the stride-2 form has a leading error factor of 4× the
3-point form on uniform mode amplitude). Phase 2's STATUS line
already records this for the scalar case
\passthrough{\lstinline!divergence(gradient(\(\psi\)))!} versus
\passthrough{\lstinline!laplacian(\(\psi\))!}.

For the manufactured field
\(v = (\sin x \cos y,\, \sin y \cos z,\, \sin z \cos x)\), the
analytical componentwise Laplacian is \(\nabla^2 v = -2\,v\): each
component depends on exactly two axes, contributing
\(-\sin(\cdot)\cos(\cdot)\) per axis when differentiated twice. The
test asserts the rate of agreement between
\passthrough{\lstinline!divergence(gradient(v))!} and the analytical
\(-2v\) on \(N \in \{16, 32, 64\}\) at slope in
\([1.7, 2.3]\).

\subsection{3. Algorithm}\label{phase13-algorithm}

\subsubsection{File layout}\label{phase13-files}

\begin{itemize}
\tightlist
\item
  \href{../../../include/gridcalc/core/EigenAliases.hpp}{\passthrough{\lstinline!include/gridcalc/core/EigenAliases.hpp!}}
  — adds \passthrough{\lstinline!Mat3d!} alongside
  \passthrough{\lstinline!Vec3d!}.
\item
  \href{../../../include/gridcalc/diff/Gradient.hpp}{\passthrough{\lstinline!include/gridcalc/diff/Gradient.hpp!}}
  — extended with the new
  \passthrough{\lstinline!Field<Vec3d> -> Field<Mat3d>!} overload.
\item
  \href{../../../include/gridcalc/diff/Divergence.hpp}{\passthrough{\lstinline!include/gridcalc/diff/Divergence.hpp!}}
  — extended with the new
  \passthrough{\lstinline!Field<Mat3d> -> Field<Vec3d>!} overload.
\item
  \href{../../../include/gridcalc/tensor/Contraction.hpp}{\passthrough{\lstinline!include/gridcalc/tensor/Contraction.hpp!}}
  — new file: \passthrough{\lstinline!trace!},
  \passthrough{\lstinline!singleContract!},
  \passthrough{\lstinline!doubleContract!}.
\item
  \href{../../../test/vector_tensor_test.cpp}{\passthrough{\lstinline!test/vector\_tensor\_test.cpp!}}
  — new file: 16 tests covering the new types, the convergence rates,
  the contraction primitives, and the three identity acceptance
  checks.
\item
  \href{../../../test/field_test.cpp}{\passthrough{\lstinline!test/field\_test.cpp!}}
  — extended with four
  \passthrough{\lstinline!FieldMat3dTest!} cases.
\end{itemize}

\subsubsection{Loop structure of the rank-2
gradient}\label{phase13-grad-loop}

\begin{lstlisting}[language={C++}]
for (int k = 0; k < grid.getNz(); ++k)
  for (int j = 0; j < grid.getNy(); ++j)
    for (int i = 0; i < grid.getNx(); ++i) {
      Mat3d M;
      for (int comp = 0; comp < 3; ++comp) {
        // Three independent stencil sums, one per axis.
        double dx = 0, dy = 0, dz = 0;
        for (m in stencil) {
          dx += w[m] * field(i + off[m], j, k)(comp);
          dy += w[m] * field(i, j + off[m], k)(comp);
          dz += w[m] * field(i, j, k + off[m])(comp);
        }
        M(comp, 0) = dx * inv_hx;
        M(comp, 1) = dy * inv_hy;
        M(comp, 2) = dz * inv_hz;
      }
      result(i, j, k) = M;
    }
\end{lstlisting}

The inner-most stencil sum is unchanged from Phase 2's scalar
gradient; the outer loop replicates it three times (once per vector
component) and packs the results into rows of the output matrix. The
\passthrough{\lstinline!Vec3d!}-valued field access
\passthrough{\lstinline!field(i + off, j, k)(comp)!} reads one
\passthrough{\lstinline!Vec3d!} (24 bytes) and extracts a single
\passthrough{\lstinline!double!} from it; modern compilers should
fold this into a single load.

The rank-2 divergence has the dual structure: one stencil sum per
output component, summed across the three axes.

\subsubsection{Contraction primitives: pointwise loops over
\passthrough{\lstinline!data()!}}\label{phase13-contraction}

\passthrough{\lstinline!trace!},
\passthrough{\lstinline!singleContract!}, and
\passthrough{\lstinline!doubleContract!} all use a single flat loop
over \passthrough{\lstinline!Field::data()!} (i-fastest order is
preserved end-to-end) and delegate to
\passthrough{\lstinline!Eigen::Matrix3d::trace()!},
\passthrough{\lstinline!operator*!}, and
\passthrough{\lstinline!(A.array() * B.array()).sum()!}
respectively. Per Phase 13's design decision, the operations are
fully materialized — every call allocates a fresh
\passthrough{\lstinline!Field!} for the output. Phase 15 may replace
these with expression-template variants alongside the
\passthrough{\lstinline!func::evaluate!}-over-tensor-fields surface
that motivates them; at Phase 13's rank-2-only scope the rank-6
explosion the ETs are meant to fend off does not yet appear.

\subsection{4. Design decisions}\label{phase13-decisions}

These distill the choices recorded in
\passthrough{\lstinline!specs/2026-05-04-phase-13-vector-tensor-fields/requirements.md!}.

\begin{itemize}
\item
  \textbf{Rank-2 type backed by \passthrough{\lstinline!Eigen::Matrix3d!}.}
  Standard Eigen module, not the unsupported tensor module — the
  Phase 0 dependency surface is unchanged. The
  \passthrough{\lstinline!TensorFixedSize!} portability risk recorded
  under \passthrough{\lstinline!STATUS.md!}'s ``Open / deferred items''
  stays deferred to Phase 15. No
  \passthrough{\lstinline!Tensor3!} /
  \passthrough{\lstinline!Tensor4!} aliases are introduced
  speculatively.
\item
  \textbf{Index convention
  \(M(i, j) = \partial_j v_i\).} Velocity-gradient form: row index =
  vector component, column index = derivative axis. Trace gives the
  divergence; the matched
  \passthrough{\lstinline!divergence(M)\(_i\) = \(\partial_j M(i, j)\)!}
  makes the third roadmap-acceptance identity drop out by
  construction. Documented on the new gradient overload's Doxygen
  block.
\item
  \textbf{Function overloading on input type.} Same
  \passthrough{\lstinline!gradient!} /
  \passthrough{\lstinline!divergence!} names; new overloads on
  \passthrough{\lstinline!Field<Vec3d>!} /
  \passthrough{\lstinline!Field<Mat3d>!}. ADL resolves the right one;
  Phase 15 may add a unifying templated form when rank-3 inputs
  arrive. The Phase 7
  \passthrough{\lstinline!Order!} parameter is preserved on both new
  overloads.
\item
  \textbf{Materialized contractions, ETs deferred to Phase 15.}
  Mirrors the Phase 11 deferral rationale. Rank-2-only inputs do not
  trigger the rank-6 explosion that motivates ETs; the materialized
  forms are simple, debuggable, and adequate for the Phase 13 test
  surface and Phase 15's eventual
  \passthrough{\lstinline!½ \(\sigma\):\(\varepsilon\)!} elastic-energy
  worked example.
\item
  \textbf{Three-identity acceptance bar instead of explicit
  \passthrough{\lstinline!curl!}.} The literal roadmap
  example \(\nabla\times\nabla\phi = 0\) is captured as Hessian
  symmetry; an explicit \passthrough{\lstinline!curl!} (rank-1
  \(\to\) rank-1) is not introduced. Adds the
  \passthrough{\lstinline!trace!}-=-divergence and
  \passthrough{\lstinline!div(grad(v))!}-=-vector-Laplacian
  identities for breadth — three independent strands of evidence the
  operators agree as a coherent algebraic system.
\end{itemize}

\subsection{5. References}\label{phase13-references}

External:

\begin{itemize}
\tightlist
\item
  \textbf{Gurtin, M. E. (1981). \emph{An Introduction to Continuum
  Mechanics}. Academic Press.} ISBN 978-0-12-309750-7.
  Section 4 (``Material and Spatial Descriptions'') and Section 6
  (``Tensor Fields'') give the velocity-gradient convention
  \((\nabla v)_{ij} = \partial_j v_i\) used throughout this phase
  and the corresponding rank-lowering divergence and Laplacian
  identities.
\item
  \textbf{Holzapfel, G. A. (2000). \emph{Nonlinear Solid Mechanics:
  A Continuum Approach for Engineering}. Wiley.} ISBN
  978-0-471-82319-3. Sections 1.4--1.6 cover the same conventions
  with explicit index-notation examples that match the Phase 13 unit
  tests.
\item
  \textbf{LeVeque, R. J. (2007). \emph{Finite Difference Methods for
  Ordinary and Partial Differential Equations}. SIAM.}
  \url{https://doi.org/10.1137/1.9780898717839}. Section 2.18 covers
  composed-operator stencil arithmetic and the order-of-magnitude
  difference between
  \(D_x^2\) and the composed
  \(D_x \circ D_x\) (the stride-2 phenomenon used in
  \Cref{phase13-stride2}).
\item
  \textbf{Eigen project: \emph{Matrix and vector types}.} Permanent
  URL:
  \url{https://eigen.tuxfamily.org/dox/group__TutorialMatrixClass.html}.
  Documents \passthrough{\lstinline!Eigen::Matrix3d!} and the
  default-constructor behavior cited in
  \passthrough{\lstinline!core/EigenAliases.hpp!}'s Doxygen block on
  \passthrough{\lstinline!Mat3d!}.
\end{itemize}

Internal cross-references (do not satisfy the ``external'' minimum):

\begin{itemize}
\tightlist
\item
  \Cref{chap:dev:gradient-and-divergence} — Phase 2 scalar
  \passthrough{\lstinline!gradient!} /
  \passthrough{\lstinline!divergence!} surface and the original
  stride-2-stencil note.
\item
  \Cref{chap:dev:higher-order-accuracy-stencils} — Phase 7
  \passthrough{\lstinline!Order!} template parameter inherited by the
  new overloads.
\item
  \Cref{chap:dev:higher-order-functional} — Phase 11 deferral
  rationale that this phase mirrors for the contraction layer.
\end{itemize}
